% Supplementary material of the manuscript (compile on its own: Overleaf Menu > Main document = supplementary.tex, pdfLaTeX + BibTeX).
% Tables and figures come from the same tables/ and figures/ folders as main.tex; numbers are never retyped.
\documentclass[preprint,12pt,authoryear]{elsarticle}
\usepackage[T1]{fontenc}
\usepackage[utf8]{inputenc}
\usepackage{lmodern}
\usepackage{microtype}
\usepackage{amsmath}
\usepackage{graphicx}
\usepackage{booktabs}
\usepackage{adjustbox}
\usepackage{xcolor}
\usepackage{siunitx}
\usepackage{setspace}
\usepackage[section]{placeins}
\usepackage{url}
\usepackage[hidelinks]{hyperref}
\onehalfspacing
\setlength{\emergencystretch}{3em}
\graphicspath{{figures/}}
\newcommand{\safefig}[2][\linewidth]{\IfFileExists{#2}{\includegraphics[width=#1]{#2}}{\fbox{\parbox[c][2.4cm][c]{0.92\linewidth}{\centering missing figure \detokenize{#2}}}}}
\newcommand{\fittable}[1]{\begin{adjustbox}{max width=\linewidth,max totalheight=0.8\textheight}\input{#1}\end{adjustbox}}
\renewcommand{\thesection}{S\arabic{section}}
\renewcommand{\thetable}{S\arabic{table}}
\renewcommand{\thefigure}{S\arabic{figure}}
\providecommand{\doi}[1]{\url{https://doi.org/#1}}

\begin{document}
\begin{center}
{\Large\bfseries Supplementary material}\\[0.6em]
{\large Explainable and hallucination-aware cross-modality learning from chest CT for X-ray disease classification: lesion-faithful bone
suppression and Swin transformer knowledge transfer}
\end{center}

\section{Deviations from the analysis plans}
\label{sup:dev}
The deviations log of the project records every change made after a plan was frozen, with its date and reason. For Study I the main
deviations are described in the main text where they arise: the encoder replacement rule and tie-break (Experiments), the JSRT overlap map
(Counterfactual lesion audit) and the post-freeze corrections (Limitations). For Study II:
\begin{itemize}
\item \emph{Gradient accumulation.} Swin-V2-B did not fit a 16\,GB GPU at batch size 32 and was screened with four micro-batches per step
  (same effective batch). The other encoders ran unchanged.
\item \emph{Encoder learning rate in the screen.} ConvNeXt-V2-T stayed at chance at the plan's learning rate. The three encoders without
  batch normalisation were screened again with the encoder learning rate reduced fourfold; this was decided before the other screen results
  were seen, on validation data only.
\item \emph{Holm families.} Before any test evaluation the authors fixed how H6 and H7 would be corrected if two encoders went forward, and
  which models H10 uses. With one encoder the families are those of the main text.
\item \emph{Execution.} Data loading, reruns of unfinished jobs and resumption across sessions changed without changing any analysis.
\item \emph{Repeated NIH runs.} One account's NIH soft-tissue array differed at the byte level from the other five copies. The five runs
  trained on it (soft seeds 0, 3, 4; early seeds 2, 3) were repeated on a consistent copy before the analysis. With the superseded runs,
  the early and soft seed ensembles would have scored 0.8279 and 0.8253 instead of 0.8273 and 0.8243; this check lies outside the plan and
  entered no test.
\end{itemize}
For Study III:
\begin{itemize}
\item \emph{CT micro-batches.} Gradient accumulation was added to the CT stage after the freeze so that four volumes of 24 slices fit a T4
  GPU; each step still averages the same four volumes, and the network has no batch statistics.
\item \emph{Teacher seed 2.} The teacher of seed 2 collapsed (constant logits; no code or data fault was found). Seeds 5 and 6 were trained
  with the unchanged recipe under a replacement rule fixed before any P1 model of seed 2 was trained; seed 5 took slot 2 (Table~\ref{tab:xm_ct_seeds}).
\item \emph{Data handling.} The design kept everything derived from CT-RATE in the cloud environment except metrics and heat-map frames, but
  its own account plan required the teacher outputs (logits and pooled features of the subset volumes, no images) on other accounts, which
  was possible only through a download; the teachers' predictions and weights and the weights of the distilled projection models also
  reached the authors' computer. Nothing derived from CT-RATE was released.
\item \emph{Execution.} One account replaced another and prepared its own NIH copy, checked file by file against the common copy; DenseNet-121
  reference runs were copied to an account with the common copy, because the rib-emphasis test reads the soft-tissue array; a rerun that
  sat in the cloud queue was resubmitted unchanged. VinBigData fine-tuning, planned only if the compute quota allowed, was not run.
\item \emph{Descriptive analyses.} The grouping of the per-class DeLong tests, the conformal analysis and the summary of the pointing game
  were fixed after the hypothesis results were known.
\end{itemize}

\section{Additional results of Study II}
\label{sup:cls}
\begin{figure}[htbp]
\centering
\safefig{figures/fig_seed_spread_densenet121.pdf}
\caption{Test macro-AUROC of every run (open circles; BDCXR-3257: each run on its own test fold) and of the seed ensembles that the tests
compare (diamonds).}
\label{fig:cls_seeds}
\end{figure}

\begin{figure}[htbp]
\centering
\safefig{figures/fig_conformal_sets_densenet121.pdf}
\caption{Conformal prediction at nominal 0.90: (a) coverage per true class and (b) mean set size of the Mondrian sets on BDCXR-3257;
(c) test false-negative rate and (d) mean set size under false-negative risk control on NIH ChestX-ray14. Dashed: the nominal level.}
\label{fig:cls_conformal}
\end{figure}

\begin{figure}[htbp]
\centering
\safefig{figures/fig_xai_ig_densenet121.pdf}
\caption{Integrated Gradients per input stream of one hybrid model (seed 0) on the NIH ChestX-ray14 test images of the main-text attribution
figure.}
\label{fig:cls_ig}
\end{figure}

\section{Additional results of Study III}
\label{sup:xm}
\begin{table}[htbp]
\centering\small
\caption{Every CT teacher run (unchanged recipe, 20 epochs): best validation and test macro-AUROC over the six findings of the CT subset.
Slot $k$ is the teacher of P1 seed $k$.}
\label{tab:xm_ct_seeds}
\fittable{tables/v4_ct_seeds.tex}
\end{table}

\begin{table}[htbp]
\centering\small
\caption{Test macro-AUROC of every single model (one seed each); the main tables use the seed ensembles. Projection models and zero-shot
transfer: six findings; NIH classifiers: 14 diseases.}
\label{tab:xm_seeds}
\fittable{tables/v4_seeds.tex}
\end{table}

\begin{table}[htbp]
\centering\footnotesize
\caption{Shared findings on the NIH test split: AUROC of the seed ensembles and paired DeLong differences with Holm over the 15 tests
(descriptive).}
\label{tab:xm_shared}
\fittable{tables/v4_shared_delong.tex}
\end{table}

\begin{table}[htbp]
\centering\small
\caption{Mean Pearson correlation of the projected CT map with the X-ray map of the same scan (matched) and of the other scans positive for
the same finding (mismatched), maps averaged over the five seed pairs.}
\label{tab:xm_crossmodal}
\fittable{tables/v4_crossmodal.tex}
\end{table}

\begin{table}[htbp]
\centering\small
\caption{Conformal risk control of the per-image false-negative rate at 0.10 on the NIH test split and selective classification by
$1-$entropy (descriptive).}
\label{tab:xm_conformal}
\fittable{tables/v4_conformal.tex}
\end{table}

\begin{table}[htbp]
\centering\footnotesize
\caption{NIH per-class test AUROC of the seed ensembles, with the paired DeLong difference F$_\mathrm{in}$ minus DenseNet-121 and Holm over
the 14 classes (descriptive).}
\label{tab:xm_per_class}
\fittable{tables/v4_per_class.tex}
\end{table}

\begin{table}[htbp]
\centering\small
\caption{Counterfactual nodules: mean change of the Mass/Nodule score after insertion, by nodule-size tertile.}
\label{tab:xm_nodule}
\fittable{tables/v4_nodule_size.tex}
\end{table}

\begin{table}[htbp]
\centering\small
\caption{Hallucination tests on the 9,861 finding-free NIH test images: share with some class above its threshold on the clean image, share
flipped by each perturbation, inserted-nodule sensitivity and share of Grad-CAM++ mass outside the lung fields for positive lung-finding
predictions.}
\label{tab:xm_halluc}
\fittable{tables/v4_halluc.tex}
\end{table}

\begin{figure}[htbp]
\centering
\safefig[88mm]{figures/fig_v4_ct_training.pdf}
\caption{Validation macro-AUROC of the CT teacher per epoch for all seven seeds. Seed 2 never left chance level, seed 6 learned late and was
not used, and seed 0 stalled near 0.60 for 12 epochs before it learned.}
\label{fig:xm_ct_training}
\end{figure}

\begin{figure}[htbp]
\centering
\safefig{figures/fig_v4_forest.pdf}
\caption{Pre-specified contrasts of Study III with 95\,\% bootstrap intervals: (a) AUROC contrasts of H11--H13; (b) H14; (c) H15. Filled
markers: $p<0.05$ (Holm-adjusted for H11 and H12).}
\label{fig:xm_forest}
\end{figure}

\begin{figure}[htbp]
\centering
\safefig{figures/fig_v4_transfer.pdf}
\caption{Macro-AUROC of single seeds (open circles) and seed ensembles (diamonds): (a) projection test split; (b, c) zero-shot transfer to
NIH and VinBigData; (d) NIH classifiers over the 14 diseases.}
\label{fig:xm_transfer}
\end{figure}

\begin{figure}[htbp]
\centering
\safefig[88mm]{figures/fig_v4_crossmodal.pdf}
\caption{Mean correlation of the CT map with the X-ray map of the same scan (filled) and of other scans positive for the same finding (open),
for P1 and P0; in parentheses the number of maps.}
\label{fig:xm_crossmodal}
\end{figure}

\begin{figure}[htbp]
\centering
\safefig{figures/fig_v4_pointing.pdf}
\caption{Pointing game: (a) hit rate over all 984 boxes with 95\,\% image-cluster bootstrap intervals; (b) focal (filled) and other (open)
boxes.}
\label{fig:xm_pointing}
\end{figure}

\begin{figure}[htbp]
\centering
\safefig{figures/fig_v4_conformal.pdf}
\caption{(a) Risk-coverage curves of the Swin classifiers (Brier score of the images kept, most confident first). (b) Test false-negative rate
under conformal risk control at 0.10 (dashed); numbers: mean label-set size.}
\label{fig:xm_conformal}
\end{figure}

\section{Additional results of Study I}
\label{sup:sup1}
\begin{table}[htbp]
\centering\small
\caption{First analysis plan, H1 (v1 family): lesion-aware arms against the recipe arm on synthetic LIDC test CTs with exact truth (three
seeds each), and the 300k-sample finals as a replication. Holm over the seven contrasts of the v1 ablation.}
\label{tab:v1_h1}
\fittable{tables/v1_h1.tex}
\end{table}

\begin{figure}[htbp]
\centering
\safefig{figures/fig_bone_uncertainty.pdf}
\caption{Exploratory recalibration on real pairs (post hoc). (a) Per-image coverage on the 51 held-out SZCH-X-Rays pairs for global and
Mondrian calibration, with quantiles pooled over pixels or taken per image; dashed: nominal 0.90. (b) Reliability of the Mondrian intervals
across nominal levels; cross: the synthetic calibration. (c) Sparsification of the absolute bone error by the ensemble spread, with the
error-sorted oracle.}
\label{fig:bone_recal}
\end{figure}

\section{Design of Study II}
\label{sup:cls_design}
\begin{figure}[htbp]
\centering
\safefig{figures/fig_study_design.pdf}
\caption{Classifier study design. (a) The three inputs for an NIH ChestX-ray14 test image: the radiograph $X$, the soft-tissue image $S$
predicted by the suppressor and the bone image $\hat B=X-S$ (box: expert annotation of a focal finding). (b) The five configurations, all with
the same DenseNet-121 encoder, class-query decoder and training settings. In the hybrid and the control one encoder with shared weights
processes each stream and the decoder attends to the tokens of all three; the control receives $X$ three times. Bottom: the pre-specified
contrasts.}
\label{fig:cls_design}
\end{figure}

\section{Per-class comparison of Study III}
\label{sup:xm_per_class}
\begin{figure}[htbp]
\centering
\safefig{figures/fig_v4_per_class.pdf}
\caption{(a) NIH per-class test AUROC of the seed ensembles of Study III. (b) F$_\mathrm{in}$ minus DenseNet-121 per class with 95\,\% DeLong
intervals; filled markers: Holm $p<0.05$ over the 14 classes. Hernia carries almost all of the H13 difference.}
\label{fig:xm_per_class}
\end{figure}

\section{Uncertainty in Study II}
\label{sup:cls_selective}
\begin{table}[htbp]
\centering\small
\caption{Conformal prediction at nominal 0.90: Mondrian sets on BDCXR-3257 (coverage, mean set size of 3 classes) and false-negative risk
control on NIH (test false-negative rate, mean set size of 14 labels).}
\label{tab:cls_unc}
\fittable{tables/v3_conformal.tex}
\end{table}

\begin{table}[htbp]
\centering\footnotesize
\caption{Selective classification: area under the risk--coverage curve (AURC; lower is better) with score A (one minus the ensemble
entropy), score A on the images that have a spread value, and score B (entropy and suppressor spread). On NIH, 16 of the 25,596 test images
have an empty lung-field mask and hence no spread, so scores A and B are compared on the same 25,580 images.}
\label{tab:cls_selective}
\fittable{tables/v3_selective.tex}
\end{table}
\begin{figure}[htbp]
\centering
\safefig{figures/fig_selective_classification_densenet121.pdf}
\caption{Risk--coverage curves for (a) BDCXR-3257 (top-1 error) and (b) NIH ChestX-ray14 (Brier score), keeping the most confident images
first. Solid: score A; dashed: score B with the suppressor spread. Legend: AURC of score A.}
\label{fig:cls_selective}
\end{figure}
\end{document}
